\chapter{讨论}\label{chap:discussion}

\section{为何正样本优先规则更具预算效率}

根据本文统计结果，正样本优先规则在与平衡规则相同的缺失率水平下获得了更低的平均采样成本。这一现象的本质在于：随着模型能力增强，简单 prompt 更容易在前几轮便产生正确答案，而继续追求“必须同时看到错误样本”的目标，将把大量预算消耗在少数 hard negative 的搜索上。对训练系统而言，这部分预算未必比转移给中等难度 prompt 更具边际收益。

因此，正样本优先规则并不是否定负样本的重要性，而是在预算有限的训练现实中重新排序不同 prompt 的优先级。它体现的是一种预算再分配思想：对简单题适度早停，对难题保留更大的探索空间。

若将这一解释与第三章中的等待时间分析对应起来，就可以更清楚地理解第四章实验结果。对平衡规则而言，有效信号出现概率由 $q_{\mathrm{bal}}(p_x;n)=1-p_x^n-(1-p_x)^n$ 决定；当 $p_x$ 逐渐接近 $1$ 时，等待一个错误样本的代价会迅速上升。第四章表~\ref{tab:ruleA} 与表~\ref{tab:ruleB} 的差异，正是这一机制的经验体现：在相同或更低的信号缺失率下，正样本优先规则显著降低了平均采样数，说明后期训练中的主要瓶颈已经不再是“找不到正确答案”，而是“为了维持平衡而额外寻找错误答案”。

同样地，训练动态分析也为这一判断提供了额外支撑。表~\ref{tab:training_dynamic_full} 显示，随着训练推进，极难样本占比下降而中高正确率区间持续扩张，这意味着越来越多 prompt 进入“正样本易得、负样本稀缺”的区域。于是，正样本优先规则的优势并不是某个孤立超参数设置下的偶然结果，而是由训练分布右移所诱发的阶段性结构变化。换言之，退出规则的最优性并非静态常数，而与模型当前能力水平和 prompt 难度分布共同决定。

从停止时间视角看，平衡规则同时受到“右端找错难”和“左端找对难”两类成本的拖累；而正样本优先规则主动放弃了前者，因此把高成本区域压缩到了低成功率一侧。只要训练中后期已经出现可观比例的高正确率 prompt，这种规则差异就会立刻转化为平均采样成本差异。换句话说，正样本优先的优势并不依赖某一个特定数据集，而是来源于其对后训练分布右移趋势的结构性适配。

\section{几何块策略的联合收益}

从纯统计角度看，几何块自适应与固定块自适应在多数设置下非常接近；但从训练系统角度看，几何块策略具有更规整的批次形态。前两轮采用小块快速筛查，后续大块则可更好地匹配现代推理引擎对同前缀多回答生成的优化需求\citep{xu2026kvcacheoptimizationstrategies,jaillet2025online,dong2025accelerating}。这意味着几何块方法有望将“统计效率接近不变”的优势转化为“系统吞吐进一步提升”的优势。

更具体地说，几何块策略的价值并不只是“把 32 次采样拆成若干轮”，而是把 prompt 的生命周期分成了两个截然不同的阶段：早期筛查阶段用极小块快速识别简单题和极难题，中后期扩展阶段则只对仍然模糊的 prompt 投放大块预算。这种结构一方面降低了无效采样占比，另一方面使活跃 prompt 集合在时间上逐渐收缩，从而为系统调度带来更高的可预测性。与随机长度或完全自由追加的采样方式相比，几何块方法更容易被实现为批量推理服务中的固定调度模板。

这一点也可以与第四章中的结果形成直接对应。尽管表~\ref{tab:ruleA} 与表~\ref{tab:ruleB} 中，几何块与固定块的信号缺失率几乎一致，但几何块在规则 B 下仍将平均采样数进一步降低到了 12.8240。换言之，它并没有通过牺牲统计质量来换取系统友好性，而是在统计效果近似不变的前提下，为批次组织提供了更规整的执行接口。因此，几何块策略的价值应被理解为“把统计上的自适应性转译成工程上的可调度性”，这也是本文强调其联合收益而非单点收益的原因。

这一点对于真实训练系统尤为关键。许多统计上“略优”的采样器在部署时并不一定高效，因为它们可能引入过于碎片化的调度、频繁的前缀失配或不可预测的批次形态。相比之下，几何块策略虽然不是最自由的在线策略，却在灵活性与可执行性之间取得了更稳健的平衡：前几轮用统一小块做宽筛，后几轮再用较大块做集中扩展，既保留了自适应能力，又保留了系统实现所需的规则性。

\section{从采样策略到隐式课程学习}

本文的另一个启示在于，自适应采样实际上会诱导一种隐式课程学习。传统课程学习通常显式定义“先学简单样本还是先学困难样本”的顺序，而本文所讨论的预算分配并不直接改变样本进入训练集的先后，而是改变不同样本在单位时间内能够获得多少 rollout 预算。其结果是：稳定全错和稳定全对的样本都会更快退出活跃集合，中等难度、结果最不确定的 prompt 会在训练过程中获得更高的持续关注度。

这种机制与显式难度排序不同。它不是在数据层面对样本做一次性划分，而是在训练过程中根据当前策略的表现反复更新样本的“活跃程度”。于是，课程不再是预先写死的，而是由模型当前能力状态和退出规则共同生成的。对 reasoning 训练来说，这种动态课程往往比静态课程更合理，因为同一道题在训练早期和后期对应的有效信息密度并不相同。

\section{与既有研究结论的关系}

本文的观察与两类既有结论形成互补。其一，RL 主要激发已有能力而非创造新能力上限的观点，强调了潜在能力暴露的重要性\citep{kim2025reinforcement}。本文则进一步指出，若采样预算设计不当，模型即使已经具备潜在正确解，也可能因为固定小样本而无法在训练时显式暴露出来。其二，自一致性与 pass@$k$ 研究强调多次采样的价值\citep{chen2021codex,wang2022selfconsistency}；本文则补充说明，关键并不只是“采得更多”，而是“把额外预算投给最需要它的 prompt”。

从这个意义上看，本文也为“后训练主要在重排已有能力”这一判断提供了更细化的机制解释：当模型已有能力以高方差形式隐含存在时，采样预算的作用并不是创造新能力，而是让训练过程更高概率地观测到那些原本会被固定小样本掩盖的正负差异。因此，预算分配可以被理解为一种能力显化机制，而不仅仅是推理成本控制机制。

进一步地，这一解释也有助于理解为什么本文特别强调 prompt 级难度异质性。若所有样本都处在近似相同的成功率区间，那么固定预算虽然未必最优，但其代价也不会太高；真正导致预算浪费急剧放大的，是同一训练切片中同时存在极难、过渡和极易三类 prompt。当这种异质性与固定小组采样相遇时，训练系统就会在两端都付出高昂的“无差异观察成本”，而自适应预算恰恰是在削减这部分成本。

\section{局限性}

本文仍存在三方面局限。第一，当前定量分析主要基于 \texttt{Qwen2.5-Math-1.5B base} 在 \texttt{OpenR1-Math-220k} 上的单模型切片数据，尚未扩展到更大模型和多基准端到端训练。第二，关于温度缩放与重要性采样的讨论目前主要停留在理论层面，尚缺少配套的完整系统实验。第三，本文从统计视角推断了几何块策略对 KV Cache 友好的潜在收益，但尚未直接报告 wall-clock 时间与吞吐测量结果。

这些局限并不削弱本文的核心结论：在 reasoning 强化学习后训练中，采样预算本身就是训练信号设计的一部分。后续工作可在更多模型、更多推理引擎和真实在线训练配置上进一步验证这一结论。

除此之外，本文还存在一个较为实际的限制：当前分析主要依据 prompt 终局正确率构造统计指标，而没有直接纳入答案长度、推理轨迹多样性以及步骤级 verifier 的反馈。如果未来训练系统更强调过程监督而非终局监督，那么“什么样的 rollout 才算形成有效信号”这一标准可能需要重新定义，相应的退出规则与预算分配器也应做相应调整。

\section{效度威胁分析}

除了一般性的局限外，本文还需要从研究效度角度对结论边界进行说明。

在内部效度方面，本文所使用的 prompt 级正确率均由有限次重复采样估计而来，尽管 $256$ 次已能提供较稳定的经验分布，但它终究只是对真实 $p_x$ 的近似。当某些样本位于极低或极高成功率区域时，有限样本误差仍可能影响全对/全错概率的估计精度。此外，本文对预算效率的讨论以二值奖励和 prompt 级终局正确率为主，没有直接建模步骤级奖励或长度惩罚，因此结论主要适用于可验证终局奖励场景。

在外部效度方面，当前数据来自单一模型与单一训练切片。虽然 OpenR1-Math-220k 具有一定代表性，但其难度构成、题目风格和采样温度设置仍可能与其他数学基准、代码推理任务或多模态 reasoning 数据存在差异。因此，本文更适合被理解为“在当前 reasoning 训练形态下发现的稳定统计规律”，而不是对所有后训练任务的无条件结论。

在构念效度方面，本文将“有效训练信号”主要定义为 prompt 内部是否同时存在可形成对比的正负样本。这一定义对于组相对归一化更新是自然的，但对于完全不同的目标函数或值函数估计方法，signal collapse 的具体表现形式可能有所变化。换言之，本文分析的是一类与组采样直接相关的训练障碍，而不是所有强化学习训练障碍的统一解释。

尽管存在这些威胁，本文的主要结论仍具有稳健性：固定小组采样在难度异质数据上会系统性地产生零信号区域，而自适应预算分配能够以更高概率恢复这些区域中的可学习差异。未来工作可以通过更大规模端到端训练、更多模型家族以及更细粒度奖励设计，对这一结论进行进一步外推与修正。

\section{实践启示与设计建议}

基于全文分析，本文给出以下几条面向实际训练系统的设计建议。

第一，对于早期试验阶段，不宜一开始就盲目增大组大小。固定大组虽然能够降低坍塌率，但其代价是线性上升的推理成本。更合理的做法是先采用小块自适应采样观察 prompt 分布，再根据分布形态决定是否值得扩大预算上限。换言之，组大小不应被视为静态超参数，而应被视为与数据难度共同决定的动态配置。

第二，在模型已经具备一定基础能力的中后期训练中，建议优先采用正样本优先退出规则。本文的统计表明，此时零信号的一大来源已经从“全错”转向“全对”，继续在简单题上强求 hard negative 的边际收益会显著下降。更优策略是让简单题较早退出，把更多预算保留给中等难度样本，以维持有效的正负对照区间。

第三，自适应采样的部署应尽量与推理引擎的前缀复用机制结合。若系统能够在同一 prompt 的多轮采样之间保留 KV Cache，则几何块策略的工程收益将比本文静态统计所体现的收益更大。这意味着未来的训练框架设计不应将“统计采样器”和“推理执行器”分离实现，而应尽量构造统一的数据流与缓存策略。

第四，训练监控指标也需要相应升级。除平均奖励、loss 和 pass@$k$ 外，prompt 级全错率、全对率、活跃 prompt 比例以及不同难度区间的预算占比，都应作为常规诊断指标纳入日志系统。原因在于，只有直接观察这些统计量，才能判断训练是否陷入了“看似稳定、实则缺少对比信号”的隐性停滞状态。

总的来说，本文更倾向于将自适应采样理解为一种训练操作系统：它不仅决定样本如何被看见，也决定预算如何被消耗，以及哪些 prompt 会在训练过程中持续保有影响力。若这一层设计得当，后续关于奖励函数、长度控制和探索温度的许多优化才更可能真正生效。
